Die Schulbehörden einer Provinz möchten einen Service anbieten,
über den Schulkinder individuelle Rechenaufgaben zum Üben generiert
bekommen können.
Wegen eines verbreiteten Aberglaubens sollen die Rechenaufgaben
aber niemals ein Resultat haben, welches eine Unglückszahl ist.

Da es sich dabei um eine relativ einfache Programmieraufgabe handelt,
glauben die Behörden zwei Fliegen mit einer Klappe schlagen zu können,
indem Sie die Programmierung eines solchen Dienstes als Informatik-Challenge 
an den Mittelschulen der Provinz ausschreiben.
Etwas überrascht vom vielfältigen Echo auf die Ausschreibung suchen sie
jetzt nach einem Tool, mit dem Sie die Eingaben filtern können, so dass
nur noch Programme genauer untersucht und bewertet werden müssen, welche
nur korrekte Aufgaben ohne Unglückszahlresultate produzieren.
Wie stehen die Chancen, ein solches Tool zu finden?

\begin{loesung}
Es geht um die Eigenschaft \textit{LUCKY} einer Sprache.
Eine Sprache $L$ hat diese Eigenschaft, wenn $L$ nur aus syntaktisch
korrekten arithmetischen Ausdrücken besteht, deren Wert verschieden
ist von $4$ (die vor allem in Japan als Unglückszahl gilt).
Die Sprache
\[
\begin{aligned}
L_1&=\{\texttt{1+3}\}&&\text{hat die Eigenschaft \textit{LUCKY} nicht,}\\
L_2&=\{\texttt{3+3+3}\}&&\text{hat die Eigenschaft \textit{LUCKY}.}
\end{aligned}
\]
Die beiden Sprachen sind von einem Programm erkennbar.
Somit ist \textit{LUCKY} eine nichttriviale Eigenschaft.
Nach dem Satz von Rice ist $\textit{LUCKY}_{\text{TM}}$ nicht
entscheidbar, ein Tool der verlangten Art kann es daher nicht geben.
\end{loesung}

\begin{diskussion}
Die Tetraphobie ist in Asien weit verbreitet, in Japan z.~B.~weil das
Kanji-Zeichen
\begin{CJK}{UTF8}{min}四\end{CJK}
für \textit{vier} und das Zeichen
\begin{CJK}{UTF8}{min}死\end{CJK}
für \textit{Tod}
in der On-Lesung gleich als
\textesh i (phonetisch) oder 
\begin{CJK}{UTF8}{min}し\end{CJK} (Hiragana) 
ausgesprochen werden.
Die Unglückszahl war in der Aufgabe nicht spezifiziert, daher wurden 
in der Prüfung viele verschiedene Unglückszahlen gewählt:
\begin{center}
\begin{tabular}{lr}
\hline
Unglückszahl       & Anzahl\raisebox{2pt}{\strut}\\
\hline
4                  &      3\raisebox{2pt}{\strut}\\
7                  &      3\\
13                 &     29\\
666                &      2\\
nicht spezifiziert &     11\\
\hline
\end{tabular}
\end{center}
Die Statistik zeigt die kulturelle Dominanz der Unglückszahl 13
in Westeuropa.
\end{diskussion}

\begin{bewertung}
Eigenschaft \textit{LUCKY} ({\bf P}) 1 Punkt,
zwei Sprachen ({\bf L}) je 1 Punkt,
Eigenschaft ist nicht trivial ({\bf T}) 1 Punkt,
Anwendung des Satzes von Rice ({\bf R}) 2 Punkte.
\end{bewertung}


